\documentclass[12pt]{article}

% ---- Кодировки и язык ----
\usepackage[utf8]{inputenc}
\usepackage[T2A]{fontenc}
\usepackage[english, main=russian]{babel}

% ---- Математика ----
\usepackage{amsmath}
\usepackage{amssymb}

% ---- Графика ----
\usepackage{graphicx}
\usepackage{subcaption}
\usepackage{float}

% ---- Таблицы ----
\usepackage{booktabs}
\usepackage{array}
\usepackage{longtable}

% ---- Интервалы и отступы ----
\usepackage{setspace}
\onehalfspacing

% ---- Поля по ГОСТ 7.32 ----
\usepackage[top=2cm, bottom=2cm, left=3cm, right=1.5cm]{geometry}
\setlength{\parindent}{1.25cm}

% ---- Гиперссылки (без цветных рамок) ----
\usepackage[hidelinks]{hyperref}

% ---- URL ----
\usepackage{url}

% ---- Заголовки разделов ----
\usepackage{textcase}
\usepackage{titlesec}

% Главы (section) — по центру, жирный, прописные
\titleformat{\section}
  {\bfseries\Large\filcenter}{\thesection.}{4pt}{\MakeTextUppercase}

% Подразделы — жирный, слева
\titleformat{\subsection}
  {\bfseries\large}{\thesubsection.}{4pt}{}

\titleformat{\subsubsection}
  {\bfseries\normalsize}{\thesubsubsection.}{4pt}{}

% ---- Подписи рисунков и таблиц по ГОСТ ----
\usepackage{caption}
\DeclareCaptionLabelFormat{gostfigure}{Рисунок #2}
\DeclareCaptionLabelFormat{gosttable}{Таблица #2}
\DeclareCaptionLabelSeparator{gost}{~---~}
\captionsetup{labelsep=gost}
\captionsetup*[figure]{labelformat=gostfigure}
\captionsetup*[table]{labelformat=gosttable}
\renewcommand{\thesubfigure}{\asbuk{subfigure}}

% ---- Листинги кода ----
\usepackage{listings}
\usepackage{xcolor}

\lstset{
  language=C++,
  basicstyle=\small\ttfamily,
  keywordstyle=\color{blue}\bfseries,
  commentstyle=\color{gray}\itshape,
  stringstyle=\color{red!70!black},
  numbers=left,
  numberstyle=\tiny\color{gray},
  numbersep=5pt,
  breaklines=true,
  frame=single,
  captionpos=b,
  inputencoding=utf8,
  extendedchars=true,
  literate={а}{{\cyra}}1 {б}{{\cyrb}}1 {в}{{\cyrv}}1 {г}{{\cyrg}}1
           {д}{{\cyrd}}1 {е}{{\cyre}}1 {ж}{{\cyrzh}}1 {з}{{\cyrz}}1
           {и}{{\cyri}}1 {й}{{\cyrishrt}}1 {к}{{\cyrk}}1 {л}{{\cyrl}}1
           {м}{{\cyrm}}1 {н}{{\cyrn}}1 {о}{{\cyro}}1 {п}{{\cyrp}}1
           {р}{{\cyrr}}1 {с}{{\cyrs}}1 {т}{{\cyrt}}1 {у}{{\cyru}}1
           {ф}{{\cyrf}}1 {х}{{\cyrh}}1 {ц}{{\cyrc}}1 {ч}{{\cyrch}}1
           {ш}{{\cyrsh}}1 {щ}{{\cyrshch}}1 {ъ}{{\cyrhrdsn}}1
           {ы}{{\cyrery}}1 {ь}{{\cyrsftsn}}1 {э}{{\cyrerev}}1
           {ю}{{\cyryu}}1 {я}{{\cyrya}}1
           {А}{{\CYRA}}1 {Б}{{\CYRB}}1 {В}{{\CYRV}}1 {Г}{{\CYRG}}1
           {Д}{{\CYRD}}1 {Е}{{\CYRE}}1 {Ж}{{\CYRZH}}1 {З}{{\CYRZ}}1
           {И}{{\CYRI}}1 {Й}{{\CYRISHRT}}1 {К}{{\CYRK}}1 {Л}{{\CYRL}}1
           {М}{{\CYRM}}1 {Н}{{\CYRN}}1 {О}{{\CYRO}}1 {П}{{\CYRP}}1
           {Р}{{\CYRR}}1 {С}{{\CYRS}}1 {Т}{{\CYRT}}1 {У}{{\CYRU}}1
           {Ф}{{\CYRF}}1 {Х}{{\CYRH}}1 {Ц}{{\CYRC}}1 {Ч}{{\CYRCH}}1
           {Ш}{{\CYRSH}}1 {Щ}{{\CYRSHCH}}1 {Ъ}{{\CYRHRDSN}}1
           {Ы}{{\CYRERY}}1 {Ь}{{\CYRSFTSN}}1 {Э}{{\CYREREV}}1
           {Ю}{{\CYRYU}}1 {Я}{{\CYRYA}}1
}

% ---- Нумерация библиографии по ГОСТ ----
\makeatletter
\renewcommand\@biblabel[1]{#1.}
\makeatother

% ---- Нумерация страниц ----
\usepackage{fancyhdr}
\pagestyle{fancy}
\fancyhf{}
\fancyhead[R]{\thepage}
\renewcommand{\headrulewidth}{0pt}

% ============================================================
\begin{document}
% ============================================================

% ============================================================
%  ТИТУЛЬНЫЙ ЛИСТ
% ============================================================
\begin{titlepage}
  \begin{center}
    \begin{scshape}
      Министерство науки и высшего образования Российской Федерации\\
      \smallskip
      Федеральное государственное автономное образовательное учреждение
      высшего образования\\
      <<Национальный исследовательский ядерный университет <<МИФИ>>
    \end{scshape}

    \vspace{1\baselineskip}
    \begin{flushleft}
      УДК 539.1
    \end{flushleft}

    \Large
    \vspace{2\baselineskip}

    \MakeTextUppercase{Выпускная квалификационная работа}\\[0.5\baselineskip]
    \begin{bfseries}
      \MakeTextUppercase{Моделирование электрических полей, дрейфа электронов
      и световыхода электролюминесценции в детекторе РЭД-100}
    \end{bfseries}

    \normalsize
    \vspace{2\baselineskip}

    \begin{minipage}[b][2\baselineskip][b]{0.4\linewidth}
      Научный руководитель\\
      (к.\,ф.-м.\,н.)
    \end{minipage}
    \hfill
    \begin{minipage}[b][2\baselineskip][b]{0.4\linewidth}
      \underline{\hspace{2cm}}~Г.\,Д. Долганов
    \end{minipage}

    \begin{minipage}[b][2\baselineskip][b]{0.4\linewidth}
      Научный консультант\\
      (Старший преподаватель)
    \end{minipage}
    \hfill
    \begin{minipage}[b][2\baselineskip][b]{0.4\linewidth}
      \underline{\hspace{2cm}}~И.\,Н. Мачулин
    \end{minipage}

    \begin{minipage}[b][2\baselineskip][b]{0.4\linewidth}
      Студент
    \end{minipage}
    \hfill
    \begin{minipage}[b][2\baselineskip][b]{0.4\linewidth}
      \underline{\hspace{2cm}}~Д.\,А. Кораблев
    \end{minipage}

    \vfill
    Москва 2026
  \end{center}
\end{titlepage}

% ============================================================
%  СОДЕРЖАНИЕ
% ============================================================
\newpage
\tableofcontents
\setcounter{page}{3}

% ============================================================
%  ПЕРЕЧЕНЬ СОКРАЩЕНИЙ И ОБОЗНАЧЕНИЙ
% ============================================================
\newpage
\section*{Перечень сокращений и обозначений}
\addcontentsline{toc}{section}{Перечень сокращений и обозначений}

\begin{tabular}{ll}
  CE$\nu$NS & когерентное упругое рассеяние нейтрино на ядре \\
            & (Coherent Elastic Neutrino-Nucleus Scattering) \\
  ЭЛ        & электролюминесценция \\
  LAr       & жидкий аргон (Liquid Argon) \\
  GAr       & газообразный аргон (Gaseous Argon) \\
  ФЭУ       & фотоэлектронный умножитель \\
  S1        & первичный сцинтилляционный сигнал \\
  S2        & вторичный электролюминесцентный сигнал \\
  ВВЭР      & водо-водяной энергетический реактор \\
  АЭС       & атомная электростанция \\
  $E$       & напряжённость электрического поля, В/см \\
  $E_\text{th}$ & пороговое поле электролюминесценции, В/см \\
  $\alpha_\text{EL}$ & коэффициент усиления ЭЛ, фотон/(В) \\
  $N_\text{ph}$ & ожидаемое число ЭЛ-фотонов на электрон \\
  $N$       & концентрация атомов газа, см$^{-3}$ \\
  $Y/N$     & приведённый световыход, $10^{-17}$\,см$^2$ \\
  $E/N$     & приведённое поле (Таунсенд), Тд ($1\,\text{Тд} = 10^{-17}$\,В\,см$^2$) \\
\end{tabular}

% ============================================================
%  ВВЕДЕНИЕ
% ============================================================
\newpage
\section*{Введение}
\addcontentsline{toc}{section}{Введение}

В современной физике элементарных частиц одной из важнейших задач является
исследование свойств нейтрино. Особый интерес представляет процесс
когерентного упругого рассеяния нейтрино на атомных ядрах
(CE$\nu$NS)~\cite{Freedman1974} — явление, при котором нейтрино низких
энергий (до нескольких десятков МэВ) взаимодействует не с отдельными
нуклонами, а со всем ядром целиком посредством обмена $Z$-бозоном. Сечение
такого взаимодействия имеет порядок $10^{-38}$\,см$^2$ и пропорционально
квадрату числа нейтронов в ядре, что делает его на несколько порядков
больше сечения рассеяния нейтрино на электроне или отдельном нуклоне.

Впервые процесс CE$\nu$NS был экспериментально зарегистрирован в 2017 году коллаборацией COHERENT. Результаты измерений оказались в хорошем согласии с предсказаниями Стандартной модели, что позволило сузить границы параметров для гипотетических взаимодействий «новой физики», выходящих за ее рамки~\cite{VanDessel2020}. Помимо этого, детальное изучение данного процесса имеет критическое значение для экспериментов по прямому поиску тёмной материи: сигналы от когерентного рассеяния нейтрино формируют неустранимый фон, который точно имитирует ожидаемый сигнал от частиц тёмной материи.

Для регистрации CE$\nu$NS требуются детекторы с предельно низким
энергетическим порогом, поскольку ядро отдачи при таком взаимодействии
уносит лишь кинетическую энергию порядка кэВ. Одним из перспективных
решений являются двухфазные эмиссионные детекторы на благородных газах.

В таких детекторах прохождение частицы через жидкую среду приводит к формированию двух последовательных сигналов. В момент взаимодействия часть энергии расходуется на возбуждение атомов и ионизацию среды с последующей рекомбинацией, порождая короткую вспышку первичной сцинтилляции — сигнал S1. Далее электроны, образованные при ионизации, под действием электрического поля дрейфуют к поверхности раздела фаз, переходят в газовую область с повышенным полем и там производят интенсивный электролюминесцентный световой сигнал S2.

Российский эксперимент \textbf{РЭД-100}~\cite{Akimov2022_RED100} является
одним из ведущих проектов в данной области. Детектор будет развернут вблизи
активной зоны реактора ВВЭР-1000 Калининской АЭС и нацелен на измерение
CE$\nu$NS, оценку техногенного и естественного фона, а также на
нейтринный мониторинг ядерных реакторов~\cite{Skorokhvatov2003}.

Точность работы детектора напрямую определяется однородностью и величиной
электрических полей в его объёме. Неоднородности поля вблизи электродных
сеток приводят к разбросу времён дрейфа электронов и неоднородности
световыхода сигнала~S2, что ухудшает координатное разрешение и
энергетическую калибровку. Поэтому численное моделирование электрических
полей, дрейфа электронов и световыхода ЭЛ является необходимым этапом
анализа и оптимизации детектора.

\textbf{Цель работы} — разработка и применение комплексной численной модели
для исследования электрических полей, дрейфа электронов и световыхода
электролюминесценции в детекторе РЭД-100.

\textbf{Задачи работы:}
\begin{enumerate}
  \item Построить трёхмерную модель электрического поля
        и исследовать влияние конфигурации
        детектора на распределение поля.
  \item Провести моделирование дрейфа электронов в газовой фазе детектора
        с использованием пакета Garfield\texttt{++} и расширения Magboltz.
  \item Рассчитать световыход электролюминесценции для различных начальных
        координат электронов и уровней жидкого аргона над гейтом.
  \item Оценить неоднородность сигнала~S2 и её зависимость от параметров детектора.
\end{enumerate}

% ============================================================
%  ГЛАВА 1. ЭКСПЕРИМЕНТ РЭД-100
% ============================================================
\newpage
\section{Эксперимент РЭД-100}

\subsection{Конфигурация установки}

Эксперимент \textbf{РЭД-100}~\cite{Akimov2022_RED100} — двухфазный
эмиссионный детектор, нацеленный на регистрацию событий CE$\nu$NS.
Установка будет расположена в непосредственной близости от активной зоны
реактора ВВЭР-1000 на Калининской АЭС (Тверская область, Россия) на
расстоянии около 19\,м от активной зоны. Реактор является мощным
источником электронных антинейтрино, что позволяет компенсировать
чрезвычайно малое сечение взаимодействия нейтрино с веществом высоким
потоком частиц.

Для подавления внешнего фона детектор окружён пассивной защитой, состоящей из 70 см воды и 5 см меди~\cite{Akimov2022_RED100}.

Сам детектор представляет собой большой криостат (Рисунок~\ref{fig:detector}),
внутри которого находится аргон в жидком и газообразном состоянии. Масса
активного вещества составляет около 100\,кг. Рабочие условия детектора:
температура 87\,К, давление 1{,}25\,атм (950\,Торр).

Изначально эксперимент РЭД-100 работал с использованием жидкого ксенона, однако в дальнейшем рабочее вещество было заменено на жидкий аргон. Это было сделано для снижения порога регистрации и подавления фоновых событий, вызванных спонтанной эмиссией одиночных электронов. Несмотря на то что сечение когерентного рассеяния нейтрино на аргоне меньше из-за меньшего числа нейтронов в ядре, более легкое ядро аргона приобретает большую энергию отдачи при взаимодействии с частицей. Такое увеличение энерговыделения переводит искомый сигнал от CE$\nu$NS в область, где влияние паразитных одноэлектронных шумов минимизировано, что существенно повышает чувствительность установки.

\begin{figure}[H]
  \centering
  \includegraphics[width=0.85\linewidth]{Autumn_2024/Figures/detector_new.jpg}
  \caption{Конструкция детектора РЭД-100.}
  \label{fig:detector}
\end{figure}

\subsection{Принцип регистрации}

Детектор работает по следующему принципу:
\begin{enumerate}
  \item Когда частица взаимодействует с жидким аргоном, происходит вспышка
        сцинтилляционного света — первичный сигнал \textbf{S1}. Одновременно
         образуются электроны ионизации.
  \item Электроны подхватываются дрейфовым электрическим полем и движутся
        вверх к поверхности раздела жидкой и газовой фаз.
  \item Достигнув поверхности, электроны вытягиваются в газовую область, где
        ускоряются в сильном электрическом поле. При столкновениях с атомами
        газа они возбуждают их; при релаксации возбуждённые атомы испускают
        фотоны с длиной волны 128 нм. Этот процесс называется
        \textbf{электролюминесценцией} (ЭЛ) и формирует вторичный,
        значительно более мощный сигнал — \textbf{S2}.
\end{enumerate}

Оба сигнала регистрируются фотоэлектронными умножителями (ФЭУ),
расположенными в верхней и нижней частях детектора. Задержка между S1 и S2 определяет вертикальную координату события (координату~$Z$), а распределение света по ФЭУ от сигнала S2, происходящего в непосредственной близости от них, позволяет восстановить координаты в плоскости~$XY$.

Принципиально важно, что амплитуда сигнала~S2 пропорциональна числу
первичных электронов, то есть энергии, переданной частицей жидкому аргону.
Поэтому однородность световыхода ЭЛ непосредственно определяет
энергетическое разрешение детектора.

\subsection{Электродная структура}

Для формирования дрейфового поля в жидком аргоне и электролюминесцентного поля в газовой фазе внутри объёма детектора натянуто несколько сетчатых электродов: катод (внизу), гейт (вблизи поверхности раздела фаз) и анод (вверху газовой фазы). Гейт представляет собой две параллельные сетки, расположенные на расстоянии 3 мм друг от друга. Сетки изготовлены из тонких проволок, образующих шестиугольный узор с шагом ячейки 4 мм.

Основное назначение гейта заключается в
экранировании газовой области от попадания электронов от фоновых событий, например, от попадания космических мюонов в детектор. Идентифицировать такие события можно по сигналу S1.

Однородность дрейфового поля в жидком аргоне достигается с помощью формирующих поле колец, окружающих активный объем.

Помимо электродов, формирующих дрейфовое и электролюминесцентное поля, в конструкции РЭД-100 присутствуют две экранирующие заземленные сетки, расположенные рядом с верхней и нижней матрицами фотоэлектронных умножителей (ФЭУ) соответственно. Главная задача этих сеток — электростатическое экранирование фотоприемников от высоких потенциалов, подаваемых на катод, анод и гейт.

Напряжения на электродах детектора приведены в Таблице~\ref{tab:voltages}.

\begin{table}[h!]
  \centering
  \caption{Напряжения на электродах детектора РЭД-100}
  \label{tab:voltages}
  \begin{tabular}{|l|c|}
    \hline
    \textbf{Электрод} & \textbf{Напряжение, В} \\
    \hline
    Анод              & $+7000$                \\
    \hline
    Гейт             & $-5000$                \\
    \hline
    Катод             & $-13000$               \\
    \hline
  \end{tabular}
\end{table}

% ============================================================
%  ГЛАВА 2. МОДЕЛИРОВАНИЕ ЭЛЕКТРИЧЕСКОГО ПОЛЯ
% ============================================================
\newpage
\section{Моделирование электрического поля}

\subsection{Постановка задачи}

Точное знание распределения электрического поля в объёме детектора
необходимо для оценки погрешностей и проведения экспериментов с различными
конфигурациями без изменения самой установки. Моделирование детектора
РЭД-100 невозможно свести к двумерной задаче, поскольку отличительной
особенностью его конструкции является шестиугольная форма ячеек
электродных сеток. Поэтому была создана трёхмерная модель центральной
части детектора.

Основная задача моделирования состоит в том, чтобы понять конфигурацию
поля в детекторе и создать максимально гибкую параметризованную модель,
пригодную для экспериментов с различными конфигурациями установки.

\subsection{Программное обеспечение}

\subsubsection*{COMSOL Multiphysics}

\textbf{COMSOL Multiphysics} --- коммерческая многофизическая
платформа для численного моделирования, основанная на методе конечных
элементов (МКЭ). Программа позволяет решать связанные задачи из различных
областей физики: электростатики, гидродинамики, теплопереноса, механики
сплошных сред и других. В настоящей работе использовался модуль
\textit{AC/DC Module}, предназначенный для расчёта электрических и
магнитных полей в трёхмерных областях произвольной геометрии.

\paragraph{Метод конечных элементов.}
Суть МКЭ состоит в разбиении расчётной области на конечное число
элементарных ячеек (элементов) --- тетраэдров, призм или гексаэдров ---
и аппроксимации искомой функции (в данном случае электрического потенциала
$\varphi$) кусочно-полиномиальными базисными функциями на каждом элементе.

\paragraph{Генерация сетки.}
COMSOL автоматически строит неструктурированную тетраэдральную сетку. Для
корректного воспроизведения поля вблизи тонких проволок анодной сетки применялось измельчение сетки:
максимальный размер элемента вблизи проволок составлял доли миллиметра,
тогда как в объёме детектора --- до $5$~мм.

\paragraph{Граничные условия и экспорт результатов.}
На электродах задавались потенциалы, соответствующие базовой конфигурации детектора: анод $+7000$~В, затворные сетки
$-5000$~В, катод $-13000$~В. Боковые грани шестиугольной призмы
моделировались как периодические, что позволяло
ограничиться одной элементарной ячейкой сетки. После решения трёхмерное
распределение поля $\mathbf{E}(x,y,z)$ экспортировалось в текстовый файл
формата, совместимого с пакетом Garfield\texttt{++} (класс
\texttt{ComponentComsol}), для последующего использования при моделировании
дрейфа электронов.

\subsection{Геометрия модели}

Моделируемая область представляет собой трёхмерную шестиугольную призму,
охватывающую центральную активную часть детектора (Рисунок~\ref{fig:geom}). Такая модель была выбрана для наглядной визуализации электрического поля.
Геометрические параметры модели приведены в Таблице~\ref{tab:sim_params_field}. 

\begin{figure}[H]
  \centering
  \includegraphics[width=0.7\linewidth]{Autumn_2024/Figures/model_geom.jpg}
  \caption{Геометрия моделируемой области.}
  \label{fig:geom}
\end{figure}

\begin{table}[h!]
  \centering
  \caption{Параметры модели для симуляции электрического поля}
  \label{tab:sim_params_field}
  \begin{tabular}{|l|c|}
    \hline
    \textbf{Параметр}                          & \textbf{Значение} \\
    \hline
    Расстояние между гейтами                   & 3\,мм             \\
    \hline
    Расстояние от верхнего гейта до анода      & 16\,мм            \\
    \hline
    Расстояние от анода до земли               & 20\,мм            \\
    \hline
    Расстояние от нижнего гейта до катода      & 415\,мм           \\
    \hline
    Уровень жидкого аргона над верхним гейтом  & 8\,мм             \\
    \hline
    Диаметр ячейки сетки                       & 4\,мм             \\
    \hline
  \end{tabular}
\end{table}

\subsection{Улучшение расчетной сетки}

Одной из проблем моделирования поля внутри детектора стала настройка
параметров расчетной сетки. Стандартные значения, предлагаемые COMSOL, дают более грубый результат. Для решения этой проблемы была проделана работа по
улучшению расчетной сетки внутри газового и жидкого объема и на
поверхности гейтов: уменьшен минимальный размер элементарной ячейки и шаг его масштабирования, что позволило увеличить число параметров модели и убрать нежелательные эффекты вблизи малых элементов геометрии.

\begin{figure}[H]
  \centering
  \includegraphics[width=0.75\linewidth]{Autumn_2024/Figures/mesh.jpg}
  \caption{Улучшенная расчётная сетка гейта.}
  \label{fig:mesh}
\end{figure}

На Рисунке~\ref{fig:mesh} показана улучшенная конфигурация расчётной
сетки на одном из гейтов и в объёме LAr. В отличие от стандартных
параметров, количество расчётных узлов увеличено за счёт уменьшения
минимального размера звеньев и увеличения степени масштабирования.

\subsection{Результаты моделирования электрического поля}

\

\subsubsection{Распределение электрического поля}

На Рисунке~\ref{fig:ep2d} показан пример распределения электрического поля
в плоскости ZY. Видно, что поле в дрейфовой области достаточно однородно, тогда как вблизи проволок анода и гейта наблюдаются значительные искажения из-за резкого перепада
напряжённости.

\begin{figure}[H]
  \centering
  \includegraphics[width=0.7\linewidth]{Autumn_2024/Figures/gate_default.jpg}
  \caption{Пример электрического поля в детекторе (в плоскости ZY).}
  \label{fig:ep2d}
\end{figure}

Важной особенностью распределения потенциала является то, что вблизи гейта создается фокусирующее электрическое поле. Оно направляет поток дрейфующих электронов сквозь центральные области ячеек сетки, тем самым предотвращая их потерю при столкновении с проволоками электрода.

\subsubsection{Исследование работы гейта}

Как было сказано ранее, одна из главных задач гейта - экранирование газовой области от электронов, возникающих из-за высокоэнергетических мюонов. Это происходит с помощью подачи дополнительного потенциала $-400В$ на верхнюю сетку гейта (Рисунок~\ref{fig:gate_closed}). 

\begin{figure}[H]
  \centering
  \includegraphics[width=0.7\linewidth]{Autumn_2024/Figures/gate_closed.jpg}
  \caption{Конфигурация электрического поля с закрытым гейтом.}
  \label{fig:gate_closed}
\end{figure}

Однако добиться полного экранирования нежелательных электронов удается не всегда. Проблема заключается в том, что в областях с низкой напряженностью электрического поля скорость дрейфа электронов существенно снижается. Это может приводить к их локальному накоплению и формированию пространственного (объемного) заряда. Собственное электрическое поле этого скопления зарядов искажает исходную конфигурацию полей детектора, ослабляя запирающие свойства гейта и провоцируя утечку фоновых электронов в область регистрации.

Для исследования явления электрического пробоя гейта определялось пороговое напряжение на аноде, при котором возникает данный процесс (Рисунок~\ref{fig:gate_proboy}).


\begin{figure}[H]
  \centering
  \begin{subfigure}[b]{0.45\textwidth}
    \includegraphics[width=\linewidth]{Autumn_2024/gate_8k_anode.jpg}
    \caption{Анод 8000\,В}
  \end{subfigure}
  \hfill
  \begin{subfigure}[b]{0.45\textwidth}
    \includegraphics[width=\linewidth]{Autumn_2024/gate_proboy.jpg}
    \caption{Анод 8250\,В}
  \end{subfigure}
  \caption{Электрическое поле вблизи гейта для различных напряжений на аноде.}
  \label{fig:gate_proboy}
\end{figure}

\subsubsection{Эксперимент с уровнем жидкого аргона}

Первым результатом, который ожидалось получить в данной работе, был график
зависимости значения электрического поля на границе раздела фаз от уровня
жидкого аргона над верхним гейтом. Перебираемые значения уровня LAr
приведены в Таблице~\ref{tab:lar_levels}.

\begin{table}[h!]
  \centering
  \caption{Исследованные значения уровня LAr над верхним гейтом}
  \label{tab:lar_levels}
  \begin{tabular}{cc}
    \toprule
    Номер & Значение уровня LAr, мм \\
    \midrule
    1  & 8      \\
    2  & 5      \\
    3  & 3      \\
    4  & 2      \\
    5  & 1      \\
    6  & 0{,}5  \\
    7  & $-$0{,}5 \\
    8  & $-$1   \\
    9  & $-$1{,}5 \\
    10 & $-$2   \\
    \bottomrule
  \end{tabular}
\end{table}

\begin{figure}[H]
  \centering
  \includegraphics[width=1.0\linewidth]{Autumn_2024/LAr_lvl_field_exp.jpg}
  \caption{Зависимость электрического поля на границе раздела фаз от уровня
           жидкого аргона над верхним гейтом.}
  \label{fig:lar_plot}
\end{figure}

Данный график (Рисунок~\ref{fig:lar_plot}) позволяет оценить минимальный
уровень жидкого аргона, подходящий для задачи детектора РЭД-100. Из
результатов эксперимента можно сделать вывод, что располагать уровень
раздела фаз ниже верхнего гейта нельзя, так как это сильно уменьшает
значение вытягивающего поля, что приведёт к значительному уменьшению
сигнала или даже к его потере. Также, начиная с уровня 3 мм LAr над гейтом, разность между средним и минимальным значением поля почти отсутствует, что говорит о хорошей однородности электрического поля для данной конфигурации детектора. 

\subsubsection{Эксперимент с положением заземлённой поверхности}

Из анализа линий электрического поля в детекторе (Рисунок~\ref{fig:ep_lines})
видно, что они имеют зонтообразную форму и сильно отличаются по своей
длине, что совместно с неоднородностью поля вдоль этих линий приводит к
неоднородности световыхода.

\begin{figure}[H]
  \centering
  \includegraphics[width=0.75\linewidth]{Autumn_2024/Figures/EP.jpg}
  \caption{Линии электрического поля в детекторе.}
  \label{fig:ep_lines}
\end{figure}

Для исследования влияния положения заземлённой поверхности над анодом на
конфигурацию поля было проведено моделирование с перебором
значения расстояния от анода до земли приведены в Таблице~\ref{tab:ground_levels}.

\begin{table}[h!]
  \centering
  \caption{Исследованные значения расстояния от анода до заземлённой поверхности}
  \label{tab:ground_levels}
  \begin{tabular}{cc}
    \toprule
    Номер & Расстояние до земли, мм \\
    \midrule
    1 & 55 \\
    2 & 45 \\
    3 & 30 \\
    4 & 20 \\
    5 & 10 \\
    6 & 7  \\
    7 & 5  \\
    \bottomrule
  \end{tabular}
\end{table}

\begin{figure}[H]
  \centering
  \begin{subfigure}[b]{0.45\textwidth}
    \includegraphics[width=\linewidth]{Autumn_2024/Figures/5.png}
    \caption{5\,мм}
  \end{subfigure}
  \hfill
  \begin{subfigure}[b]{0.45\textwidth}
    \includegraphics[width=\linewidth]{Autumn_2024/Figures/7.png}
    \caption{7\,мм}
  \end{subfigure}
  \vspace{0.2cm}
  \begin{subfigure}[b]{0.45\textwidth}
    \includegraphics[width=\linewidth]{Autumn_2024/Figures/20.png}
    \caption{20\,мм}
  \end{subfigure}
  \hfill
  \begin{subfigure}[b]{0.45\textwidth}
    \includegraphics[width=\linewidth]{Autumn_2024/Figures/55.png}
    \caption{55\,мм}
  \end{subfigure}
  \caption{Результаты эксперимента при разных расстояниях до заземлённой поверхности.}
  \label{fig:ground_exp}
\end{figure}

Согласно результатам моделирования (Рисунок~\ref{fig:ground_exp}), уменьшение
расстояния между анодом и землёй сильно влияет на линии напряжённости вблизи
проволоки. При малых расстояниях возможные траектории электронов становятся
более схожими по своей длине, что потенциально может уменьшить неоднородность
световыхода. Однако говорить о полном устранении неоднородности сигнала мы не можем, так как электрическое поле продолжает оставаться неоднородным вдоль отдельных треков, поэтому необходимо дополнительно провести оценку световыхода для различных начальных положений электронов.

\subsection{Выводы по главе}

В рамках данного этапа работы было проведено моделирование электрического
поля в детекторе РЭД-100. Полученная модель отличается своей гибкостью
благодаря параметризации решения. Установлено, что уровень жидкого аргона
над верхним гейтом должен быть положительным, иначе вытягивающее поле
резко падает. Уменьшение расстояния от анода до заземлённой поверхности
влияет на форму линий поля вблизи проволок сетки.

% ============================================================
%  ГЛАВА 3. МОДЕЛИРОВАНИЕ ДРЕЙФА ЭЛЕКТРОНОВ
% ============================================================
\newpage
\section{Моделирование дрейфа электронов}

\subsection{Постановка задачи}

Ключевой фактор, определяющий качество работы детектора, --- это то,
насколько эффективно и точно можно собрать ионизационный заряд,
образовавшийся при взаимодействии частицы. Электроны должны продрейфовать
через часть детектора, заполненную веществом в жидкой фазе, затем выйти
в газовую область и создать световой сигнал. Любые искажения электрического
поля на этом пути могут привести к потере или неоднородности светового
сигнала или ошибкам в определении координат события.

В рамках данного этапа работы проводилось моделирование дрейфа электронов
в газовой фазе детектора для определения влияния неоднородности поля
вблизи анодной сетки.

\subsection{Программное обеспечение}

Для моделирования дрейфа электронов в газовой фазе детектора РЭД-100
использовалась связка из трёх программных инструментов: COMSOL Multiphysics
(расчёт поля), Magboltz (транспортные свойства газа) и Garfield\texttt{++}
(трекинг электронов). Ниже приведено подробное описание каждого из них.

\subsubsection*{COMSOL Multiphysics}

Трёхмерное распределение электрического поля, полученное в Главе~2,
экспортировалось из COMSOL в текстовый файл, содержащий координаты узлов
сетки и соответствующие значения компонент вектора $\mathbf{E}(x,y,z)$.
Формат файла совместим с классом \texttt{ComponentComsol} пакета
Garfield\texttt{++}, что позволяет напрямую загружать результаты
конечно-элементного расчёта без дополнительной конвертации. Интерполяция
поля в произвольной точке пространства выполняется методом барицентрических
координат внутри тетраэдрального элемента сетки.

\subsubsection*{Magboltz}

\textbf{Magboltz}~\cite{magboltz} --- программа для расчёта транспортных
свойств электронов в газовых смесях методом Монте-Карло путём численного
решения уравнения Больцмана для функции распределения электронов по
скоростям. Программа разработана С.~Биги (S.~Biagi) и широко применяется
в физике детекторов частиц.

\paragraph{Принцип работы.}
Magboltz моделирует движение ансамбля электронов в однородном электрическом
поле с учётом упругих и неупругих столкновений с атомами газа. Для каждого
типа столкновения используются дифференциальные сечения, хранящиеся во
встроенной базе данных. В результате расчёта определяются:
\begin{itemize}
  \item скорость дрейфа $v_d(E)$ --- средняя скорость направленного движения
        электрона вдоль поля;
  \item продольный $D_L$ и поперечный $D_T$ коэффициенты диффузии;
  \item коэффициент Таунсенда $\alpha$ (ударная ионизация);
  \item коэффициент прилипания $\eta$ (захват электрона молекулами газа).
\end{itemize}

\paragraph{Использование в работе.}
В настоящей работе Magboltz запускался для чистого аргона при температуре
$T = 87$~К и давлении $P = 950$~Торр (1,25~атм). Расчёт выполнялся для
набора значений приведённого поля $E/N$ в диапазоне, охватывающем все
значения поля, встречающиеся в объёме детектора. Полученная таблица
$v_d(E)$ загружалась в Garfield\texttt{++} через класс \texttt{MediumMagboltz}
и использовалась при интегрировании уравнений движения электрона.

\subsubsection*{Garfield\texttt{++}}

\textbf{Garfield\texttt{++}}~\cite{garfield} --- объектно-ориентированный
программный пакет для детального моделирования детекторов частиц,
регистрирующих ионизацию в газовых и жидких средах. Пакет написан на языке
C++ и является развитием оригинального кода Garfield (FORTRAN, CERN),
переработанного с учётом современных требований к архитектуре программного
обеспечения. Garfield\texttt{++} интегрирован с ROOT (CERN) и Magboltz.

\paragraph{Архитектура пакета.}
Основные классы Garfield\texttt{++}, использованные в работе:
\begin{itemize}
  \item \texttt{ComponentComsol} --- загрузка и интерполяция поля из файла
        COMSOL; обеспечивает метод \texttt{ElectricField(x,y,z,...)} для
        получения вектора $\mathbf{E}$ в произвольной точке;
  \item \texttt{MediumMagboltz} --- интерфейс к Magboltz; хранит таблицу
        транспортных свойств газа и предоставляет метод
        \texttt{ElectronVelocity(ex,ey,ez,bx,by,bz,vx,vy,vz)} для
        вычисления вектора скорости дрейфа при заданном поле;
  \item \texttt{DriftLineRKF} --- трекинг электрона методом Рунге--Кутты
        с адаптивным шагом по времени.
\end{itemize}

\paragraph{Алгоритм трекинга.}
Движение электрона описывается уравнением
\begin{equation}
  \frac{d\mathbf{r}}{dt} = \mathbf{v}_d\!\left(\mathbf{E}(\mathbf{r})\right),
\end{equation}
где $\mathbf{v}_d$ --- вектор скорости дрейфа, определяемый из таблицы
Magboltz по текущему значению поля $|\mathbf{E}|$. Интегрирование
выполняется методом Рунге--Кутты 4-го порядка с автоматическим выбором
шага $\Delta t$, обеспечивающим заданную точность позиции. На каждом шаге
вычисляется вклад в световыход электролюминесценции согласно модели
Бузулуцкова (см. Главу~4). Трекинг завершается при достижении электроном
анода или выходе за пределы расчётной области.

\subsection{Выбор геометрии}

Электроды в РЭД-100 --- это не сплошные пластины, а сетки с шестиугольными
звеньями. Расчёт поля сразу во всём детекторе занимает много времени.
Поэтому было выбрано одно повторяющееся звено сетки с условиями
периодичности на боковых гранях. Стоит отметить, что изначальная модель создавалась для наглядной визуализации конфигурации электрического поля в детекторе, поэтому она не оптимальна с точки зрения моделирования дрейфа. 

\begin{figure}[H]
  \centering
  \begin{subfigure}{0.48\textwidth}
    \centering
    \includegraphics[width=\linewidth, height=12cm]{Autumn_2024/Figures/field_init.jpg}
    \caption{}
    \label{fig:field_full}
  \end{subfigure}
  \hfill
  \begin{subfigure}{0.48\textwidth}
    \centering
    \includegraphics[width=\linewidth, height=12cm]{Autumn_2024/Figures/field_optimized.jpg}
    \caption{}
    \label{fig:field_opt}
  \end{subfigure}
  \caption{Электрическое поле вблизи сетки: (а)~изначальная конфигурация;
           (б)~оптимизированная конфигурация с условием периодичности.}
  \label{fig:field_comparison}
\end{figure}

На Рисунке~\ref{fig:field_comparison} можно заметить, что поле в
изначальной модели совпадает с полем в оптимизированной модели с условием
периодичности. Это позволило создать геометрию с меньшим числом расчётных
параметров, но с той же точностью моделирования.

\begin{figure}[H]
  \centering
  \includegraphics[width=0.9\linewidth]{Autumn_2024/Figures/geom_optimized.jpg}
  \caption{Геометрия моделируемой области (одна шестиугольная ячейка).}
  \label{fig:geom_drift}
\end{figure}

\subsection{Транспортные свойства электронов}

Зная поле, можно задать движение электронов в объёме детектора. Параметры
дрейфа электрона были получены с помощью расширения \textbf{Magboltz} для
инструмента Garfield\texttt{++}. Расчёт выполнялся при рабочих условиях
детектора: температура 87\,К, давление 950\,Торр (1{,}25\,атм).

\begin{figure}[H]
  \centering
  \includegraphics[width=0.9\linewidth]{Autumn_2024/Figures/velocity.png}
  \caption{Скорость дрейфа электронов в газовом аргоне как функция
           напряжённости поля (Magboltz, 87\,К, 950\,Торр).}
  \label{fig:velocity}
\end{figure}

Полученная таблица скоростей использовалась Garfield\texttt{++} для
вычисления локальной скорости электрона в каждой точке траектории.
Объединение значений электрического поля из COMSOL и физики газа из
Magboltz внутри Garfield\texttt{++} позволило получить треки дрейфующих
электронов.

\subsection{Результаты моделирования дрейфа}

В связи с тем, что полученная модель детектора оказалась компактной,
удалось набрать десятки тысяч треков электронов.

\subsubsection{Анализ времени дрейфа}

Одним из важных параметров, которые нужно было оценить, являлось среднее
время дрейфа электрона в газовой фазе:
\begin{itemize}
  \item \textbf{Среднее время дрейфа:} 1633{,}31\,нс;
  \item \textbf{Медианное вермя дрейфа:} 1607{,}00\,нс.
\end{itemize}

Данные значения были получены с помощью запуска электронов, распределённых
равномерно по поверхности раздела фаз.

\begin{figure}[H]
  \centering
  \includegraphics[width=0.9\linewidth]{Autumn_2024/drift_time_cropped.jpg}
  \caption{Гистограмма распределения времени дрейфа электронов.}
  \label{fig:drift_hist}
\end{figure}

На Рисунке~\ref{fig:drift_hist} можно подробнее рассмотреть распределение
времён дрейфов электронов. На гистограмме справа можно заметить тяжёлый
хвост --- это электроны, запущенные из центра ячейки. Данная область
наиболее интересна, так как, как было показано ранее, основная часть электронов в реальном
эксперименте проходит именно через центр шестиугольной ячейки из-за фокусировки при прохождении через гейт.

\subsubsection{Радиальная зависимость времени дрейфа}

\begin{figure}[H]
  \centering
  \includegraphics[width=0.9\linewidth]{Autumn_2024/Figures/time_vs_r.jpg}
  \caption{Распределение времени дрейфа от расстояния до центра ячейки.}
  \label{fig:radial_drift}
\end{figure}

Для наглядного понимания распределения времени дрейфа от расстояния до
центра ячейки был построен график на Рисунке~\ref{fig:radial_drift}. По
нему можно доказать, что наибольшее время дрейфа имеют электроны,
запущенные из центра. Также было произведено фитирование данного
распределения с помощью экспоненты. Значение $R^2 = 0{,}9828$, что близко
к единице и говорит о том, что независимые параметры фитирования с высокой
точностью описывают полученное распределение.

\subsubsection{Траектории дрейфа}

\begin{figure}[H]
  \centering
  \includegraphics[width=0.9\linewidth]{Autumn_2024/Figures/tracks.jpg}
  \caption{Траектории дрейфа электронов в разных проекциях.}
  \label{fig:tracks}
\end{figure}

На Рисунке~\ref{fig:tracks} показаны траектории дрейфа электронов,
запущенных из центра со случайным отклонением в плоскости XY на 1\,мм.
По графикам видно, что они имеют зонтообразную форму, что совпадает с
ожидаемым результатом, полученным при моделировании в COMSOL
(Рисунок~\ref{fig:ep_lines}).

\subsection{Выводы по главе}

Проведённое моделирование дрейфа электронов в газовой фазе детектора
РЭД-100 позволило установить следующее:
\begin{itemize}
  \item Среднее время дрейфа электрона в газовой фазе составляет около
        1633\,нс (при запуске электронов равномерно по всей границе раздела фаз)
  \item Время дрейфа монотонно убывает с ростом расстояния от центра
        ячейки до точки старта электрона, что объясняется зонтообразной
        формой траекторий.
  \item Электроны, стартующие из центра ячейки, имеют наибольшее время
        дрейфа и проходят наибольший путь в газовой фазе.
\end{itemize}

% ============================================================
%  ГЛАВА 4. МОДЕЛИРОВАНИЕ СВЕТОВЫХОДА ЭЛЕКТРОЛЮМИНЕСЦЕНЦИИ
% ============================================================
\newpage
\section{Моделирование световыхода электролюминесценции}

\subsection{Постановка задачи}

Принципиально важно, что сигнал~S2 пропорционален числу первичных
электронов, то есть энергии, переданной частицей жидкому аргону. Именно
поэтому амплитуда S2 непосредственно определяет энергетическое разрешение
детектора: чем больше фотонов удаётся собрать с каждого электрона, тем
точнее можно восстановить энергию исходного события.

В идеале число фотонов не должно зависеть от координаты вылета первичного
электрона из жидкой области детектора, однако траектория движения электрона
может повлиять на финальный световыход. Именно эта неоднородность
исследуется в данной главе.

\subsection{Модель электролюминесценции}

\subsubsection{Универсальная газовая формулировка}

В физике газовых разрядов выход ЭЛ описывается как функция приведённого
поля $E/N$ (в единицах Таунсенда, $1\,\text{Тд} = 10^{-17}$\,В\,см$^2$).
Согласно литературным данным~\cite{Oliveira2011}, абсолютная зависимость
для чистого газообразного аргона имеет вид:

\begin{equation}
  \frac{Y}{N} = 0{,}081 \frac{E}{N} - 0{,}190
  \quad \text{(в единицах } 10^{-17}\,\text{см}^2\text{)},
  \label{eq:EL_universal}
\end{equation}

где $Y$ --- количество фотонов, рождённых электроном на 1\,см пути,
$N$ --- концентрация среды (см$^{-3}$). Стоит отметить, что данное приближение работает в ограниченной области напряжений (Рисунок~\ref{fig:tracks}).

\begin{figure}[H]
  \centering
  \includegraphics[width=0.9\linewidth]{Autumn_2024/argon_modelling_graph.jpg}
  \caption{Экспериментальные зависимости приведенного световыхода от значения приведенного поля для комнатной температуры~\cite{Buzulutskov2012}}.
  \label{fig:tracks}
\end{figure}

\subsubsection{Переход к пошаговому интегрированию}

Умножим уравнение~(\ref{eq:EL_universal}) на $N$ и сгруппируем термины,
извлекая константу выхода $\alpha_\text{EL}$ и порог $E_\text{th}$:

\begin{equation}
  Y = 0{,}081 \cdot E - 0{,}190 \cdot 10^{-17} \cdot N
    = 0{,}081 \left( E - 2{,}3457 \cdot 10^{-17} \cdot N \right).
  \label{eq:EL_linear}
\end{equation}

Таким образом, ожидаемое число фотонов, испускаемых электроном на
элементарном отрезке пути~$dl$, описывается формулой
Бузулуцкова~\cite{Buzulutskov2012}:

\begin{equation}
  \frac{dN_\text{ph}}{dl} = \alpha_\text{EL} \cdot \max\bigl(0,\, E(l) - E_\text{th}\bigr),
  \label{eq:EL}
\end{equation}

где $E(l)$~[В/см] --- напряжённость поля вдоль траектории,
$E_\text{th}$~[В/см] --- пороговое поле электролюминесценции,
$\alpha_\text{EL} = 0{,}081$~[В$^{-1}$] --- коэффициент усиления.
Данное приближение характерно для электрических полей порядка~кВ/см;
при дальнейшем увеличении поля основной вклад начинает вносить ионизация.

Интегрирование по полной траектории даёт суммарный световыход:

\begin{equation}
  N_\text{ph} = \int_0^{s_\text{tot}} \alpha_\text{EL} \cdot
    \max\bigl(0,\, E(s) - E_\text{th}\bigr)\, ds.
  \label{eq:EL_total}
\end{equation}

\subsubsection{Применение к условиям детектора}

Для условий детектора РЭД-100 (1{,}25\,атм и 87\,К) концентрация атомов
по уравнению состояния идеального газа $N = p / (k_B T)$ составляет:

\begin{equation}
  N = \frac{1{,}25 \cdot 101325}{1{,}38 \cdot 10^{-23} \cdot 87}
    \approx 1{,}054 \cdot 10^{20}\,\text{см}^{-3}.
  \label{eq:N}
\end{equation}

Соответствующее пороговое поле электролюминесценции:

\begin{equation}
  E_\text{th} = 2{,}3457 \cdot 10^{-17} \cdot 1{,}054 \cdot 10^{20}
    \approx 2473\,\text{В/см}.
  \label{eq:Eth}
\end{equation}

Финальный алгоритм расчёта на каждом шаге интегрирования:

\begin{equation}
  dN_\text{ph} = 0{,}081 \cdot \max(0,\, E - 2473) \cdot dl,
  \label{eq:EL_step}
\end{equation}

где $E$~[В/см] --- напряжённость поля в текущей точке, $dl$~[см] --- длина
шага.

\subsection{Профиль поля вдоль оси детектора}

Для понимания того, в какой области траектории происходит основной вклад
в световыход, был построен профиль напряжённости электрического поля
$E(z)$ вдоль различных траекторий электронов (Рисунок~\ref{fig:field_profile}).

\begin{figure}[H]
  \centering
  \includegraphics[width=0.9\linewidth]{Autumn_2024/E_field_vs_t.jpg}
  \caption{Профиль напряжённости электрического поля $E(z)$ вдоль оси
           детектора (уровень LAr~5\,мм). Горизонтальная пунктирная линия
           соответствует порогу электролюминесценции
           $E_\text{th} \approx 2473$\,В/см. Область выше порога
           вносит вклад в световыход~S2.}
  \label{fig:field_profile}
\end{figure}

Из Рисунка~\ref{fig:field_profile} видно, что электрон, запущенный из центра ячейки, проходит через область пониженного электрического поля, а для других электронов видно, что электрическое поле вдоль всей траектории имеет значение выше порогового. 

\subsection{Процедура трекинга}

На каждом шаге интегрирования $\Delta t = 0{,}1$\,нс выполнялись следующие
операции:
\begin{enumerate}
  \item Вычисление напряжённости поля $\mathbf{E}$ в текущей точке через
        интерполяцию по сетке COMSOL.
  \item Определение скорости дрейфа $\mathbf{v}$ из таблицы Magboltz.
  \item Вычисление пройденного пути $dl = |\mathbf{v}|\,\Delta t$ и
        обновление координат.
  \item Накопление числа фотонов
        $dN_\text{ph} = \alpha_\text{EL}\,\max(0, E - E_\text{th})\,dl$.
\end{enumerate}

Трекинг завершался при выходе электрона из газовой области. Для каждой
конфигурации уровня~LAr запускалось \textbf{10000 электронов}, стартовые
координаты которых равномерно распределялись по квадрату
$[-0{,}2;\,0{,}2]\times[-0{,}2;\,0{,}2]$\,см в плоскости~XY, покрывая
одну полную ячейку сетки.

\subsection{Результаты моделирования световыхода}

\subsubsection{Профиль световыхода вдоль траектории}

На Рисунке~\ref{fig:path_profile} показан кумулятивный световыход
$N_\text{ph}(s)$ и $Z(s)$ как функции пройденного пути~$s$ и вертикальной
координаты~$z$ для типичной траектории при уровне LAr~8\,мм. Вклад в
световыход начинается лишь после входа электрона в область поля выше
порогового~$E_\text{th}$: это соответствует зазору между гейтовой сеткой
и анодом, где поле наиболее велико. Участки траектории в межгейтовом
зазоре и в жидком аргоне в световыход практически не вносят.

\begin{figure}[H]
  \centering
  \includegraphics[width=0.9\linewidth]{Autumn_2024/Figures/cumul_z.jpg}
  \caption{Кумулятивный световыход $N_\text{ph}$ в зависимости от
           пройденного пути~$s$ и вертикальной координаты~$z$
           (уровень LAr~5\,мм).}
  \label{fig:path_profile}
\end{figure}

\subsubsection{Тепловая карта световыхода в плоскости раздела фаз}

На Рисунке~\ref{fig:heatmap} представлена тепловая карта ожидаемого числа
фотонов $N_\text{ph}(x_0, y_0)$ для уровня~LAr~8\,мм. Карта демонстрирует
отчётливую неоднородность: минимальный световыход наблюдается в центре
ячейки, а вблизи проволок гейтовой сетки он увеличивается. Это объясняется тем, что около электродов напряжённость поля сильно возрастает. Следовательно, электроны, чьи траектории начинаются над проволоками гейта, проходят через области с более сильным полем по сравнению с электронами, дрейфующими из центра ячейки, что математически выражается в росте интеграла~(\ref{eq:EL_total}) и, как следствие, в большем  световыходе. Кроме того, согласно Рисунку~\ref{fig:field_profile}, видно, что электроны, запущенные из центра, проходят через область пониженного электрического поля, где, согласно модели, не образуются ЭЛ-фотоны.

\begin{figure}[H]
  \centering
  \includegraphics[width=0.85\linewidth]{Autumn_2024/Figures/5mm_hm.jpg}
  \caption{Тепловая карта световыхода $N_\text{ph}(x_0, y_0)$ для уровня
           LAr~8\,мм.}
  \label{fig:heatmap}
\end{figure}

\subsubsection{Зависимость световыхода от расстояния до центра ячейки}

Также была получена зависимость количества фотонов от расстояния до центра
ячейки $N_\text{ph}(r_0)$, где $r_0 = \sqrt{x_0^2 + y_0^2}$ --- расстояние
от центра ячейки до точки старта. На Рисунке~\ref{fig:radial_photons}
показана эта зависимость для уровня LAr~5\,мм. Световыход монотонно растёт
с ростом~$r_0$: электроны из центра ячейки производят заметно меньше
фотонов, чем электроны, стартующие вблизи проволок. Величина такой неоднородности составляет примерно 10\%.

\begin{figure}[H]
  \centering
  \includegraphics[width=0.9\linewidth]{Autumn_2024/Figures/5mm_vs_r0.jpg}
  \caption{Зависимость ожидаемого числа ЭЛ-фотонов от радиуса спавна~$r_0$
           (уровень LAr~5\,мм). Полоса обозначает стандартное отклонение
           внутри бина.}
  \label{fig:radial_photons}
\end{figure}

\subsubsection{Влияние уровня жидкого аргона}

На Рисунке~\ref{fig:heatmaps_all} приведены тепловые карты $N_\text{ph}(x_0, y_0)$
для всех исследованных уровней от 2 до 12\,мм. С увеличением уровня~LAr
средний световыход уменьшается (Рисунок~\ref{fig:photons_vs_lvl}): более
толстый слой жидкости смещает поверхность раздела фаз дальше от гейтовой
сетки и делает область газового дрейфа меньше. При этом неоднородность
световыхода по ячейке при больших уровнях выражена сильнее, чем при
меньших.

\begin{figure}[H]
  \centering
  \begin{subfigure}{0.45\textwidth}
    \centering
    \includegraphics[width=\linewidth, height=7cm]{Autumn_2024/Figures/2mm_hm.jpg}
    \caption{2\,мм}
  \end{subfigure}
  \hfill
  \begin{subfigure}{0.45\textwidth}
    \centering
    \includegraphics[width=\linewidth, height=7cm]{Autumn_2024/Figures/5mm_hm.jpg}
    \caption{5\,мм}
  \end{subfigure}
  \vfill
  \begin{subfigure}{0.45\textwidth}
    \centering
    \includegraphics[width=\linewidth, height=7cm]{Autumn_2024/Figures/7mm_hm.jpg}
    \caption{7\,мм}
  \end{subfigure}
  \hfill
  \begin{subfigure}{0.45\textwidth}
    \centering
    \includegraphics[width=\linewidth, height=7cm]{Autumn_2024/Figures/9mm_hm.jpg}
    \caption{9\,мм}
  \end{subfigure}
  \vfill
  \begin{subfigure}{0.45\textwidth}
    \centering
    \includegraphics[width=\linewidth, height=7cm]{Autumn_2024/Figures/12mm_hm.jpg}
    \caption{12\,мм}
  \end{subfigure}
  \caption{Тепловые карты световыхода для уровней LAr 2, 5, 7, 9, 12\,мм
           над верхним гейтом.}
  \label{fig:heatmaps_all}
\end{figure}

\begin{figure}[H]
  \centering
  \includegraphics[width=0.9\linewidth]{Autumn_2024/Figures/lar_lvl_vs_photons.jpg}
  \caption{Зависимость ожидаемого числа ЭЛ-фотонов от уровня жидкого
           аргона над гейтом.}
  \label{fig:photons_vs_lvl}
\end{figure}

\subsubsection{Влияние напряжения на аноде}

На Рисунке~\ref{fig:photons_vs_anode} представлена зависимость ожидаемого числа фотонов электролюминесценции от напряжения на аноде. Увеличение потенциала усиливает электрическое поле в газовой области, что повышает световыход.

\begin{figure}[H]
  \centering
  \includegraphics[width=0.9\linewidth]{Autumn_2024/el_photons_vs_anode.jpg}
  \caption{Зависимость ожидаемого числа ЭЛ-фотонов от напряжения на аноде.}
  \label{fig:photons_vs_anode}
\end{figure}

С практической точки зрения максимизация напряжения необходима для повышения энергетической чувствительности детектора при поиске низкоэнергетических событий. Однако верхний предел рабочего потенциала ограничен возникновением электрического пробоя газового кармана и эффективностью работы затвора, что приводит к ионизации газовой среды и засвету ФЭУ.

\subsection{Выводы по главе}

В данной главе было проведено моделирование световыхода электролюминесценции
в газовой фазе детектора РЭД-100. Для расчёта ожидаемого числа фотонов на
каждой траектории использовалась линейная модель~(\ref{eq:EL}) с
параметрами, масштабированными на рабочие условия детектора (87\,К,
1{,}25\,атм): $\alpha_\text{EL} = 0{,}081$\,В$^{-1}$,
$E_\text{th} \approx 2473$\,В/см.

Было показано, что основной вклад в световыход S2 формируется на участке
траектории в газовом кармане между верхним гейтом и анодом. Построенная
тепловая карта $N_\text{ph}(x_0, y_0)$ демонстрирует отчётливую
неоднородность в пределах одной ячейки шестиугольной сетки: минимум
световыхода наблюдается в центре ячейки, максимум --- над проволоками.
Усреднение по азимутальному углу подтвердило, что световыход монотонно
растёт с ростом расстояния от центра ячейки до точки старта электрона.

В ходе исследования влияния уровня жидкого аргона (от 2
до 12\,мм) установлено, что с ростом уровня средний световыход уменьшается,
а неоднородность по ячейке увеличивается.

Таким образом, неоднородности поля вблизи проволок сетки приводят к
разбросу световыхода в зависимости от поперечных координат электрона, что
ухудшает координатное разрешение и энергетическую калибровку.

% ============================================================
%  ЗАКЛЮЧЕНИЕ
% ============================================================
\newpage
\section*{Заключение}
\addcontentsline{toc}{section}{Заключение}

В рамках выпускной квалификационной работы была разработана и применена
комплексная численная модель для исследования электрических полей, дрейфа
электронов и световыхода электролюминесценции в двухфазном аргоновом
детекторе РЭД-100.

\textbf{По результатам моделирования электрического поля} (Глава~2):
\begin{itemize}
  \item Построена трёхмерная конечно-элементная модель центральной части
        детектора в COMSOL Multiphysics с оптимизированной расчётной сеткой.
        Гибкая параметризация модели позволяет проводить эксперименты с
        различными конфигурациями установки без изменения самого детектора.
  \item Установлено, что уровень жидкого аргона над верхним гейтом должен
        быть положительным: расположение поверхности раздела фаз ниже
        гейта резко снижает вытягивающее поле и может привести к потере
        сигнала. Также установлено, что, начиная с уровня 3 мм над гейтом, разница между средним и минимальным значением поля снижается, то есть растет степень однородности вытягивающего поля. 
  \item Показано, что уменьшение расстояния от анода до заземлённой
        поверхности влияет на форму линий поля вблизи проволок сетки и
        потенциально может уменьшить неоднородность световыхода.
\end{itemize}

\textbf{По результатам моделирования дрейфа электронов} (Глава~3):
\begin{itemize}
  \item Среднее время дрейфа электрона в газовой фазе составляет около
        1633\,нс (для электронов запущенных равномерно по всей площади ячейки).
  \item Время дрейфа монотонно убывает с ростом расстояния от центра
        ячейки до точки старта электрона. Электроны из центра ячейки
        имеют наибольшее время дрейфа и проходят наибольший путь.
  \item Траектории электронов имеют зонтообразную форму, что согласуется
        с результатами моделирования электрического поля в COMSOL.
\end{itemize}

\textbf{По результатам моделирования световыхода ЭЛ} (Глава~4):
\begin{itemize}
  \item Разработана и реализована процедура расчёта световыхода ЭЛ по
        линейной модели Бузулуцкова с параметрами, масштабированными на
        рабочие условия детектора (87\,К, 1{,}25\,атм):
        $\alpha_\text{EL} = 0{,}081$\,В$^{-1}$,
        $E_\text{th} \approx 2473$\,В/см.
  \item Основной вклад в световыход S2 формируется на участке траектории
        в газовом кармане между верхним гейтом и анодом, где поле
        превышает порог ЭЛ.
  \item Тепловая карта $N_\text{ph}(x_0, y_0)$ демонстрирует отчётливую
        неоднородность в пределах одной ячейки шестиугольной сетки:
        минимум световыхода наблюдается в центре ячейки, максимум -
        над проволоками затвора.
  \item С ростом уровня жидкого аргона над гейтом средний световыход
        уменьшается, а неоднородность по ячейке увеличивается.
\end{itemize}

Полученные результаты позволяют количественно оценить вклад неоднородности
поля вблизи гейтовой сетки в разброс амплитуды сигнала~S2. В дальнейшем
планируется использование построенных карт световыхода для разработки
поправок к координатной зависимости сигнала при анализе экспериментальных
данных РЭД-100.

% ============================================================
%  СПИСОК ИСПОЛЬЗОВАННЫХ ИСТОЧНИКОВ
% ============================================================
\newpage
\renewcommand{\refname}{Список использованных источников}
\addcontentsline{toc}{section}{Список использованных источников}

\begin{thebibliography}{12}

\bibitem{Freedman1974}
Freedman~D.~Z.
Coherent neutrino nucleus scattering as a probe of the weak neutral current~//
Phys. Rev. D. ---
1974. ---
Vol.~9. ---
P.~1389--1392.

\bibitem{VanDessel2020}
Van~Dessel~N. et al.
Nuclear structure aspects of spin-(in)dependent coherent elastic
neutrino-nucleus scattering~//
Phys. Rev. C. ---
2020. ---
Vol.~101. ---
035502.

\bibitem{Akimov2022_RED100}
Akimov~D.~Yu. et al. (RED-100 Collaboration).
First ground-level laboratory test of the two-phase xenon emission detector
RED-100~//
JINST. ---
2022. ---
Vol.~17. ---
P11011.

\bibitem{Skorokhvatov2003}
Skorokhvatov~M.~D., Sabelnikov~A.~V.
Neutrino monitoring of nuclear reactors~//
Phys. Atom. Nucl. ---
2003. ---
Vol.~66. ---
P.~301--309.

\bibitem{Buzulutskov2012}
Buzulutskov~A.
Electroluminescence and electron avalanche in two-phase argon detectors~//
Instruments and Experimental Techniques. ---
2012. ---
Vol.~55, No.~1. ---
P.~60--80.

\bibitem{Oliveira2011}
Monteiro~C.~M.~B. et al.
Secondary scintillation yield in pure argon~//
JINST. ---
2011. ---
Vol.~6. ---
P06003.

\bibitem{magboltz}
Biagi~S.~F.
Monte Carlo simulation of electron drift and diffusion in counting gases
under the influence of electric and magnetic fields~//
Nucl. Instrum. Methods Phys. Res. A. ---
1999. ---
Vol.~421. ---
P.~234--240.

\bibitem{garfield}
Veenhof~R.
Garfield, recent developments~//
Nucl. Instrum. Methods Phys. Res. A. ---
1998. ---
Vol.~419. ---
P.~726--730.

\end{thebibliography}

\end{document}